% ==========================================================================
% INFO 4617 Web Data Science -- Week 12: Social & Media Platform APIs (Chapter 12)
% Build:  latexmk -pdf week-12.tex   (from this folder; no env vars needed)
% (or run `make` from the slides/ directory to build every deck)
% ==========================================================================
\documentclass[9pt,aspectratio=169]{beamer}
% Resolve the shared theme/preamble in ../common/ when compiling this
% file directly from its own folder (no TEXINPUTS needed).
\makeatletter\def\input@path{{../common/}}\makeatother
\input{preamble}

\title[Week 12 · Social APIs]{{\Huge Social \& Media Platform APIs}}
\subtitle{Week 12 · APIs module · Textbook Chapter 12}
\author[Keegan]{\textbf{Brian C. Keegan, Ph.D.} \\ Associate Professor, Department of Information Science \\ University of Colorado Boulder}
\date{November 2, 4 \& 6, 2026}

\begin{document}

\maketitle

% -------------------------------------------------------------------------
\begin{frame}{This week at a glance}
    \begin{columns}[T]
        \begin{column}{0.6\textwidth}
            Last week the data came from governments obliged to publish it. This week it comes from platforms that share it \textit{at their discretion}, and can stop. We read what federal agencies post, every class live coding, from one lab notebook in three parts.

            \bigskip

            \dayband{Monday}{Lab, part 1: Bluesky}
            \begin{itemize}\small
                \item An open API with no key: who posts, and is it the agency?
            \end{itemize}

            \medskip

            \dayband{Wednesday}{Lab, part 2: more of Bluesky}
            \begin{itemize}\small
                \item More posts than one request returns, and an agency's rhythm
            \end{itemize}

            \medskip

            \dayband{Friday}{Lab, part 3: YouTube; then your project}
            \begin{itemize}\small
                \item A keyed API for the same agencies' videos; 15 minutes on your portfolio checkpoint
            \end{itemize}
        \end{column}
        \begin{column}{0.35\textwidth}
            \textbf{The lab}\\[2pt]
            {\footnotesize \texttt{labs/week-12/lab-12.ipynb} in the course repository. Each part is added before its day.\par}
            \medskip
            \begin{alertblock}{One key, for Friday}
                \footnotesize \texttt{YOUTUBE\_KEY}, a YouTube Data API key from a \textbf{personal} Google account. Set it as an environment variable, close the terminal, open a new one, and start Jupyter from there.
            \end{alertblock}
            \textbf{Due Sunday, Nov 8}\\[2pt]
            {\footnotesize Chapter 12's \textbf{Recommended Exercises}, from its companion notebook, and portfolio checkpoint 2, \textit{Collect and clean your data}. Graded on good-faith attempt to complete rather than getting everything working.\par}
            \medskip
        \end{column}
    \end{columns}
\end{frame}

% =========================================================================
\section{Monday · Live-Coding Lab}
% =========================================================================

\begin{frame}{This week, every class is live coding}
    \begin{columns}[T]
        \begin{column}{0.56\textwidth}
            \dayband{Find it}{\texttt{labs/} in the course repository}
            \begin{enumerate}\small
                \item Open the \href{https://github.com/cuinfoscience/INFO4617-Fall2026/tree/main/labs}{\textbf{labs} folder} of \texttt{cuinfoscience/INFO4617-Fall2026} on GitHub
                \item Open \texttt{week-12}, click \texttt{lab-12.ipynb}, then \textbf{Download raw file}
                \item Open it in JupyterLab in the \texttt{webdata} environment
            \end{enumerate}

            \medskip
            \dayband{Use it}{one part a day}
            \begin{itemize}\small
                \item I write the code on screen, and you write it in your copy
                \item Put \textbf{your} email in the User-Agent in Step 0
                \item Today: Part 1, Steps 0--3; no key, no account
                \item Wednesday and Friday: download the new copy and start at that part's setup step
            \end{itemize}
        \end{column}
        \begin{column}{0.4\textwidth}
            \textbf{Solutions}\\[2pt]
            {\footnotesize \texttt{lab-12-key.ipynb}, beside it: every cell filled in and run. Fall behind? Copy the cell from the key and keep going.\par}
            \medskip
            \textbf{Nothing to submit}\\[2pt]
            {\footnotesize The lab is practice. After each step, a \textbf{Discuss} box holds questions we'll talk over before the next.\par}
            \medskip
            \begin{alertblock}{Error flavors}
                \footnotesize A \texttt{KeyError} on a field means the answer wasn't what you expected: print it and look. A request that fails for everyone is \textit{broken}: tell me, or open an issue.
            \end{alertblock}
        \end{column}
    \end{columns}
\end{frame}

\begin{frame}{Platforms share at their discretion}
    \begin{columns}[T]
        \begin{column}{0.56\textwidth}
            Government APIs exist because a law says the data is public. A platform's API exists because the company chooses to offer it, and it can choose again:
            \begin{itemize}\small
                \item \textbf{Reddit} closed self-service API access on November 11, 2025: a new app now needs Reddit's approval.
                \item \textbf{Spotify} has required every new developer app's owner to have Premium since February 11, 2026, and removed top tracks, popularity, and follower counts.
                \item \textbf{Bluesky} is built on an open protocol: its public API answers anyone, with no account or key.
            \end{itemize}
            \small Chapter 12 tells the Reddit and Spotify stories; the lab reads Bluesky, then YouTube.
        \end{column}
        \begin{column}{0.4\textwidth}
            \begin{alertblock}{Social data is people data}
                \footnotesize A person made almost every post. Read the terms of service; collect the minimum you need; don't quote a post if a search could find its author. This week we read \textbf{federal agencies}, and keep private people's posts and handles out of every output.
            \end{alertblock}
        \end{column}
    \end{columns}
\end{frame}

% The code on these frames is the lab's Part 1 (labs/week-12/lab-12-key.ipynb), run live on
% 2026-11-02 against Bluesky's public AppView. Chapter 12 at textbook baf0bc0 still uses atproto
% with a login; the lab follows the chapter's planned revision.
\begin{frame}[fragile]{From the website to the API: one agency's profile}
    \begin{columns}[T]
        \begin{column}{0.56\textwidth}
            The National Weather Service at \href{https://bsky.app/profile/nws.noaa.gov}{bsky.app/profile/nws.noaa.gov}, and the same profile from the API:
\begin{lstlisting}
bsky_base = "https://public.api.bsky.app/xrpc"
profile = requests.get(
    f"{bsky_base}/app.bsky.actor.getProfile",
    params={"actor": "nws.noaa.gov"},
    headers=HEADERS).json()
print(profile["handle"], "|", profile["displayName"])
print("followers:", profile["followersCount"],
      "| posts:", profile["postsCount"])
# nws.noaa.gov | National Weather Service
# followers: 63020 | posts: 1394
\end{lstlisting}
            \small The method's name is part of the address. No key, no login: Bluesky's \textbf{AppView} answers reads for anyone.
        \end{column}
        \begin{column}{0.4\textwidth}
            \textbf{Is it really the agency?}\\[2pt]
            {\footnotesize Two proofs, both in the API's answer:\par}
            \begin{itemize}\footnotesize
                \item The \textbf{handle is a domain}, \texttt{nws.noaa.gov}: only whoever controls the domain can take it
                \item A \textbf{blue check}: \texttt{"verification"} says Bluesky itself verified it
            \end{itemize}
            \medskip
            {\footnotesize A search for ``National Park Service'' also finds ``Alt National Park Service'', with neither. Collect its posts as the agency's, and your study is wrong.\par}
        \end{column}
    \end{columns}
\end{frame}

\begin{frame}[fragile]{Five agencies in one table}
    \begin{columns}[T]
        \begin{column}{0.5\textwidth}
\begin{lstlisting}
agencies = ["noaa.gov", "nws.noaa.gov",
            "climate.noaa.gov", "noaacomms.noaa.gov",
            "federalreserve.gov"]
rows = []
for handle in agencies:
    p = requests.get(
        f"{bsky_base}/app.bsky.actor.getProfile",
        params={"actor": handle}, headers=HEADERS).json()
    rows.append({"handle": p["handle"],
                 "followers": p["followersCount"],
                 "posts": p["postsCount"],
                 "joined": p["createdAt"][:10]})
accounts = pd.DataFrame(rows).sort_values(
    "followers", ascending=False)
\end{lstlisting}
        \end{column}
        \begin{column}{0.47\textwidth}
            {\footnotesize
            \begin{tabular}{@{}lrrr@{}}
                \toprule
                \textbf{Handle} & \textbf{Followers} & \textbf{Posts} & \textbf{Joined} \\
                \midrule
                nws.noaa.gov & 63,020 & 1,394 & 2025-01 \\
                climate.noaa.gov & 46,966 & 93 & 2025-01 \\
                noaa.gov & 41,937 & 150 & 2025-01 \\
                federalreserve.gov & 8,405 & 1,029 & 2024-12 \\
                noaacomms.noaa.gov & 5,620 & 117 & 2025-01 \\
                \bottomrule
            \end{tabular}\par}
            \medskip
            {\footnotesize These are the federal accounts with a \texttt{.gov} handle on November 2; four are NOAA's. NASA and the Census Bureau have none.\par}
            \medskip
            \textbf{Microsprint}\\[2pt]
            {\footnotesize Not every official account has a domain: \texttt{librarycongress.bsky.social} and \texttt{stlouisfed.bsky.social} (FRED's home) are both verified instead.\par}
        \end{column}
    \end{columns}
\end{frame}

\begin{frame}[fragile]{What an agency posted this week}
    \begin{columns}[T]
        \begin{column}{0.53\textwidth}
\begin{lstlisting}
params = {"actor": "nws.noaa.gov", "limit": 25,
          "filter": "posts_no_replies"}
feed = requests.get(
    f"{bsky_base}/app.bsky.feed.getAuthorFeed",
    params=params, headers=HEADERS).json()["feed"]
posts = []
for item in feed:
    if "reason" in item:   # a repost
        continue
    post = item["post"]
    posts.append({"time": post["indexedAt"][:16],
                  "likes": post["likeCount"],
                  "reposts": post["repostCount"],
                  "text": post["record"]["text"][:70]})
\end{lstlisting}
        \end{column}
        \begin{column}{0.44\textwidth}
            \textbf{Most reposted, NWS}\\[2pt]
            {\footnotesize
            \begin{tabular}{@{}rl@{}}
                \toprule
                \textbf{Reposts} & \textbf{Post} \\
                \midrule
                80 & Daylight Saving Time ends \ldots \\
                33 & Updated Extreme Rainfall \ldots \\
                33 & Here's a more detailed look \ldots \\
                \bottomrule
            \end{tabular}\par}
            \medskip
            \textbf{Why skip reposts?}\\[2pt]
            {\footnotesize A repost carries someone else's post, and they may be a private person. Skipping them keeps every handle in our table an agency's.\par}
            \medskip
            {\footnotesize The answer ends with a \texttt{cursor}. Wednesday: what is it for?\par}
        \end{column}
    \end{columns}
\end{frame}

\begin{frame}[fragile]{Before Friday: your YouTube key}
    \begin{columns}[T]
        \begin{column}{0.56\textwidth}
            \begin{enumerate}\small
                \item Sign in to the \href{https://console.cloud.google.com/}{Google Cloud console} with a \textbf{personal} Google account, and create a project
                \item In its \textbf{API Library}, enable \textbf{YouTube Data API v3}
                \item Under \textbf{Credentials}, create an \textbf{API key}, and restrict it to the YouTube Data API v3
                \item Store it with the \texttt{webdata} environment, in a terminal (on Windows, the Anaconda Prompt):
            \end{enumerate}
\begin{lstlisting}
conda env config vars set YOUTUBE_KEY="your-key" -n webdata
\end{lstlisting}
            \small Then close that terminal, open a new one, run \texttt{conda activate webdata}, and start Jupyter from it.
        \end{column}
        \begin{column}{0.4\textwidth}
            \begin{alertblock}{It's a password}
                \footnotesize Never paste the key into a notebook cell or a code file. If it leaks, delete it in the console and make a new one.
            \end{alertblock}
            \textbf{The quota}\\[2pt]
            {\footnotesize 10,000 units a day; most reads cost 1 unit, and searches are capped at 100 a day (\href{https://developers.google.com/youtube/v3/getting-started}{Google's guide}). Friday's code reads channels and videos, not searches.\par}
        \end{column}
    \end{columns}
\end{frame}

% =========================================================================
\section{Wednesday · Live-Coding Lab}
% =========================================================================

\begin{frame}{Today: Part 2, more of Bluesky}
    \begin{itemize}
        \item Follow the \texttt{cursor} to read more posts than one request returns, as Wikipedia's \texttt{continue} did in week 10
        \item An agency's rhythm: when it posts, and which posts travel
        \item Download the new \texttt{lab-12.ipynb}; start at Part 2's setup step
    \end{itemize}
    \medskip
    \small This part's frames are added before Wednesday's class, after chapter 12's revision.
\end{frame}

% =========================================================================
\section{Friday · Live-Coding Lab}
% =========================================================================

\begin{frame}{Today: Part 3, YouTube's Data API}
    \begin{itemize}
        \item A keyed API: \texttt{YOUTUBE\_KEY}, read with week 11's safe ways
        \item The same agencies' channels and videos, and the API's daily quota
        \item Download the new \texttt{lab-12.ipynb}; start at Part 3's setup step
    \end{itemize}
    \medskip
    \small This part's frames are added before Friday's class.
\end{frame}

\begin{frame}{Friday's last 15 minutes: portfolio checkpoint 2}
    \begin{columns}[T]
        \begin{column}{0.56\textwidth}
            \dayband{Due Sunday}{Collect and clean your data}
            \begin{itemize}\small
                \item \textbf{Collection code} that runs, makes polite requests (a User-Agent that identifies you, pauses between requests), and saves the raw data
                \item \textbf{A cleaning step}, with its code and output, that produces a tidy dataset
                \item \textbf{A small sample} of the cleaned data in your repository, with a data dictionary
                \item \textbf{The README's Data section}: the source, when you collected it, how many records, and what you kept out of the repository and why
            \end{itemize}
        \end{column}
        \begin{column}{0.4\textwidth}
            \textbf{Use the 15 minutes}\\[2pt]
            {\footnotesize Run your collection code once, here, where you can ask. A platform this week? Its terms of service decide what goes in your repository.\par}
            \medskip
            \textbf{Graded}\\[2pt]
            {\footnotesize Two criteria, one point each, on what's pushed by Sunday night.\par}
        \end{column}
    \end{columns}
\end{frame}

% =========================================================================
\begin{frame}{Key takeaways}
    \begin{enumerate}
        \item A platform's API exists \textbf{at the platform's discretion}: Reddit and Spotify have closed much of theirs since 2024. Date what you collect, and say what access you had.
        \item \textbf{Check who is posting} before you collect: a domain handle or a verified check, not a familiar name.
        \item \textbf{Social data is people data}: read public institutions' accounts, and keep private people's posts and handles out of what you save and show.
        \item \textbf{An open protocol} answers anyone; a \textbf{keyed API} counts what you use. Both still need a User-Agent.
    \end{enumerate}
\end{frame}

\end{document}
